\chapter{Conclusions and outlook}
\section{Summary}

\subsection{Physical origin of the observed shape}

The optical photon distribution is mainly determined by the interaction depth and subsequent optical transport. Low-energy gammas deposit energy near the entrance and produce a broad response. At higher energies, deposits near the readout add a sharp central peak, while the outer component remains dominant in photon count. Gamma scattering has a comparatively small effect on the complete normalized map, despite its frequent occurrence, especially at high energies. It nevertheless enhances both the innermost core, through localized recoil deposition, and the tails, through later interactions at displaced positions.

The dominant first-interaction process changes from photoelectric absorption at low energies to Compton scattering at high energies. Escape after Compton scattering becomes more frequent with increasing energy.

BGO has a larger share of photoelectric interactions and stronger reabsorption of scattered gammas. Its shorter gamma displacements make the deposited-energy distribution more compact. This does not produce a consistent optical-shape advantage because stronger attenuation near the entrance and greater survival of long optical paths increase the broad light component.

\subsection{Parametric representation}

The four-parameter Power--Moffat mixture describes the sharp central concentration and broad wings at perpendicular incidence, with median held-out TV values of 0.020 for \ce{CeBr3} and 0.018 for BGO across thirteen energies. This agreement applies at the one-millimetre fit resolution and does not guarantee an accurate sub-millimetre peak height. For inclined beams, the two-centre model achieves a median held-out TV of 0.033, but leaves an artificial minimum between the shoulder and peak where the data show a continuous ridge.

\section{Limitations}
\subsection{Simulation model}

The continuous readout omits pixelation, wavelength-dependent detection efficiency, and electronics. The optical distribution also depends on assumed optical properties, particularly the inherited \ce{LaBr3} attenuation surrogate for \ce{CeBr3} and the wavelength-independent BGO absorption length.

Constant scintillation yields exclude nonproportionality and quenching. Such effects could change the relative light contributions of different deposits and interaction histories.

\subsection{Dataset and coverage}

The results apply to one crystal geometry and selected beam conditions. Numerical widths and fitted parameters cannot be transferred directly to other detector configurations. The maximum simulated energy of \qty{1000}{\kilo\electronvolt} also excludes pair production, so its effect on containment and the optical core-to-tail balance has not been tested.

\subsection{Interpretation and validation}

The parametric models describe optical maps at specific beam energies and angles. They have not yet been used to reconstruct unknown beam parameters. Whether these descriptions also reproduce measured detector signals remains an open experimental question.

\section{Outlook}

A direct extension of the inclined model would be a continuous connection between the shoulder and peak. The existing inclined samples provide a basis for testing such a modification without changing the detector simulation. Independently of this improvement, the current models could serve as a starting point for beam reconstruction. This would first require establishing how their shape parameters depend jointly on gamma energy and incidence angle. Tests under beam conditions not used to establish this relationship could then determine which beam parameters can be recovered from spatial shape alone and where different conditions remain indistinguishable.

Further simulations could test how the observed behaviour extends beyond the present crystal geometry and energy range. Varying the crystal thickness would change both the available interaction depths and the optical propagation distances, allowing their influence on the central peak and broad component to be studied systematically. Extending the gamma energy range above the pair-production threshold would also test how the additional charged-particle tracks and annihilation photons affect the optical distribution.

Experimental validation would be needed to establish whether the simulated dependence of the optical distribution on gamma energy and incidence angle is also observed in a physical detector. A comparison with measurements could test the predicted central peak and tails and help assess the optical-property assumptions. This would require accounting for the response of the detector used in the comparison.
